\chapter{Trading Volatility}\label{ch:dv:trading-volatility}

A trader buys a one-month at-the-money straddle at 18 volatility because the share has realised 25 for two months, and delta-hedges it every
day. A month later the share has again realised exactly 25, and the position has lost money. For three weeks the share drifted up half a
percent a day, too quietly to pay the straddle's time decay while its gamma was large. In the last week it moved almost 3\% a day, but by then it
was eight percent above the strike, where the straddle had almost no gamma left. The same 25\% realised, spread evenly over the month, would have
made 3.28 on a premium of 4.15. This chapter is about what a volatility position actually earns: \omterm{def:dv:trading-volatility:scalping}{gamma scalping} and its dependence on the path,
carry and roll-down, skew and term trades, the daily attribution of a volatility book, and the premium that sellers of volatility collect.

\section{Gamma scalping}

Chapter 4 derived the \omterm{def:dv:greeks-and-the-hedging-pnl:hpnl}{hedging P\&L} of a delta-hedged option: each day it earns $\frac12\Gamma S^2\bigl(r^2-\sigma_{\text{imp}}^2\Delta t\bigr)$,
where $r$ is the day's return. Positive when the day's move beats the implied variance, negative otherwise, weighted by the \omterm{def:dv:greeks-and-the-hedging-pnl:cashgamma}{cash gamma} of that
day.

\begin{definition}[Gamma scalping]\label{def:dv:trading-volatility:scalping}
\emph{Gamma scalping}\index{gamma scalping} is holding a long option position and re-hedging its delta at intervals, so that the hedge sells
after rises and buys after falls; its P\&L is the sum over the intervals of the \omterm{def:dv:greeks-and-the-hedging-pnl:cashgamma}{cash gamma} times the difference between the squared realised
return and the implied variance of the interval.
\end{definition}

\begin{definition}[Implied--realised spread]\label{def:dv:trading-volatility:spread}
The \emph{implied--realised spread}\index{implied--realised spread} of an option position is the difference between the implied volatility at
which it was traded and the volatility subsequently realised by the underlying over its life, the quantity a delta-hedged position is a bet
on, and only on average.
\end{definition}

The sum is weighted, and the weights move with the path. A straddle's \omterm{def:dv:greeks-and-the-hedging-pnl:cashgamma}{cash gamma} is largest near the strike and near expiry, and small far from
the strike. Realised variance earned where the gamma is small counts for little.

\begin{example}[Same volatility, different days]\label{ex:dv:trading-volatility:days}
A 21-day at-the-money straddle on a share at 100, bought at 18 volatility (premium 4.15), hedged daily at 18. Three paths all realise exactly
25\%:
\begin{itemize}
\item \textbf{even}: a move of $\pm1.57\%$ every day, up and down in turn. The straddle makes 3.28;
\item \textbf{the wrong days}: fifteen days of $+0.5\%$ (7.9\% realised), which take the share to 107.8, then six days of $\pm2.84\%$ (45\%
realised). The quiet weeks lose 0.80 while the gamma is large; the wild week, far from the strike, makes 0.02. The total is $-0.78$;
\item \textbf{the right days}: fifteen days of $\pm0.5\%$ at the strike, then the same wild week there, with the most gamma of the month. The
straddle makes 2.87.
\end{itemize}
The cash-gamma sum reproduces these within a few tenths (3.05, $-0.86$, 3.03): the rest is the higher-order terms of the large moves
(\cref{fig:dv:trading-volatility:scalp}).
\end{example}

\begin{omfigure}
\begin{tikzpicture}
\begin{groupplot}[group style={group size=2 by 1,horizontal sep=1.6cm},
  width=6.6cm,height=5cm,xlabel={trading day},xmin=0,xmax=21,
  tick label style={font=\small},label style={font=\small},grid=major,grid style={gray!25},table/col sep=comma]
\nextgroupplot[ylabel={share price},ymin=94,ymax=112,
  legend columns=3,legend cell align=left,legend style={font=\footnotesize,at={(1.12,-0.3)},anchor=north,draw=none,fill=none,/tikz/every even column/.append style={column sep=8pt}}]
\addplot[omDef,thick] table[x=day,y=s_uniform]{figdata/derivatives/25-trading-volatility/scalp.csv};
\addplot[omProp,very thick] table[x=day,y=s_wrong]{figdata/derivatives/25-trading-volatility/scalp.csv};
\addplot[omThm,thick,dashed] table[x=day,y=s_right]{figdata/derivatives/25-trading-volatility/scalp.csv};
\addplot[gray,dotted,domain=0:21] {100};
\legend{even,wrong days,right days}
\nextgroupplot[ylabel={cumulative hedged P\&L},ymin=-2,ymax=4]
\addplot[omDef,thick] table[x=day,y=pnl_uniform]{figdata/derivatives/25-trading-volatility/scalp.csv};
\addplot[omProp,very thick] table[x=day,y=pnl_wrong]{figdata/derivatives/25-trading-volatility/scalp.csv};
\addplot[omThm,thick,dashed] table[x=day,y=pnl_right]{figdata/derivatives/25-trading-volatility/scalp.csv};
\end{groupplot}
\end{tikzpicture}

\omcaption{A one-month straddle bought at 18 volatility and hedged daily along three paths that all realise 25\%. Left: the share (dotted: the
strike). Right: the cumulative P\&L. On the wrong days the share leaves the strike quietly and moves wildly where the straddle has no gamma
left. Data: the tutorial.}\label{fig:dv:trading-volatility:scalp}
\end{omfigure}

Two consequences follow. A delta-hedged option is a bet on the \omterm{def:dv:trading-volatility:spread}{implied--realised spread} only on average over paths. On one path it is a bet on
the gamma-weighted realised variance, which is what a \omterm{def:dv:variance-swaps-and-volatility-derivatives:vs}{variance swap} (chapter 14) removes by holding constant \omterm{def:dv:greeks-and-the-hedging-pnl:cashgamma}{cash gamma}. And the hedging
frequency matters little to the expected P\&L, but it matters to its noise and to transaction costs, the trade-off chapter 26 prices.

\section{Volatility carry and the roll-down}

Hold an option with nothing moving and it still earns or loses: time decay, and the change in its implied volatility as its expiry shortens.

\begin{definition}[Volatility roll-down]\label{def:dv:trading-volatility:rolldown}
The \emph{volatility roll-down}\index{volatility roll-down} of an option over a horizon is the change in its implied volatility when its time to
expiry shortens by the horizon along an unchanged term structure (and, for a sticky-delta surface, an unchanged smile in moneyness); with the
time decay, it makes up the option's carry.
\end{definition}

\begin{example}[Rolling down an upward-sloping curve]\label{ex:dv:trading-volatility:carry}
The at-the-money term structure is 19.70 at one month, 20.92 at two, 21.79 at three and 23.89 at one year. A three-month straddle worth 8.69
held for a month, with spot and curve unchanged, is worth 6.81: a carry of $-1.88$. At an unchanged 21.79 it would have lost 1.59 to time
decay. The other 0.28 is roll-down, the straddle's implied volatility falling 0.87 point as it becomes a two-month option. A calendar (long
the three-month straddle, short 1.73 one-month straddles, vega-neutral) earns 0.13 a day of carry instead, paid for with a short gamma of
$-0.17$ (\cref{fig:dv:trading-volatility:carry}).
\end{example}

\begin{omfigure}
\begin{tikzpicture}
\begin{groupplot}[group style={group size=2 by 1,horizontal sep=1.6cm},
  width=6.6cm,height=5cm,
  tick label style={font=\small},label style={font=\small},grid=major,grid style={gray!25},table/col sep=comma]
\nextgroupplot[xlabel={expiry (months)},ylabel={at-the-money volatility (\%)},xmin=0,xmax=12,ymin=18,ymax=25]
\addplot[omThm,very thick] table[x=months,y=vol]{figdata/derivatives/25-trading-volatility/term.csv};
\addplot[omProp,only marks,mark=*,mark size=2pt] table[x=months,y=vol]{figdata/derivatives/25-trading-volatility/term_points.csv};
\nextgroupplot[xlabel={trading day},ylabel={straddle value},xmin=0,xmax=21,ymin=6.5,ymax=9,
  legend columns=1,legend cell align=left,legend style={font=\footnotesize,at={(0.5,-0.3)},anchor=north,draw=none,fill=none}]
\addplot[omProp,very thick] table[x=day,y=roll]{figdata/derivatives/25-trading-volatility/carry.csv};
\addplot[omDef,thick,dashed] table[x=day,y=flat]{figdata/derivatives/25-trading-volatility/carry.csv};
\legend{rolling down the curve,volatility held at 21.79}
\end{groupplot}
\end{tikzpicture}

\omcaption{Left: an upward-sloping at-the-money term structure (dots at one, two and three months). Right: a three-month straddle held for a month
with spot and curve unchanged: its value falls by time decay and by the roll-down of its volatility. Data: the
tutorial.}\label{fig:dv:trading-volatility:carry}
\end{omfigure}

\section{Skew and term trades}

A calendar is a term trade: long one expiry's volatility, short another's, vega-neutral in total. It profits if the curve flattens or inverts, and
its carry comes from the roll-down. A risk reversal is a skew trade: long a 25-delta call, short a 25-delta put (or the reverse), nearly
vega-neutral and long or short the skew. What it earns depends on how the surface moves with the spot, which is the marking rule of chapter 7.

\begin{example}[One bad day, two marks]\label{ex:dv:trading-volatility:rr}
Three months to expiry, at-the-money volatility 20\% and a skew that puts the 25-delta put (strike 93.6) at 21.3 and the 25-delta call (106.96) at
18.7. The desk is short the put, long the call, hedged with 0.5 shares. The share falls 5\% in a day and the at-the-money level does not move.
\begin{center}\small
\begin{tabular}{lrrrrr}
\toprule
marking & gamma & \omterm{def:dv:greeks-and-the-hedging-pnl:greeks}{vega} & \omterm{def:dv:greeks-and-the-hedging-pnl:greeks}{vanna} & unexplained & total\\
\midrule
\omterm{def:dv:implied-volatility-and-its-surface:sticky}{sticky strike} & 0.053 & $-0.003$ & 0.000 & $-0.077$ & $-0.023$\\
\omterm{def:dv:implied-volatility-and-its-surface:sticky}{sticky delta} & 0.053 & $-0.003$ & 0.112 & $-0.085$ & 0.081\\
\bottomrule
\end{tabular}
\end{center}
Under \omterm{def:dv:implied-volatility-and-its-surface:sticky}{sticky strike} each option keeps its volatility. Under \omterm{def:dv:implied-volatility-and-its-surface:sticky}{sticky delta} both options' volatilities fall by a point, because each strike is now
higher relative to the spot and the skew lowers volatility as moneyness rises; the position's \omterm{def:dv:greeks-and-the-hedging-pnl:greeks}{vanna} turns that into a gain. The mark decides the sign of the day. The Greek
explain misses 0.08 either way: a 5\% move is too large for second-order terms, and a desk revalues such days in full.
\end{example}

\section{Attributing a volatility book's P\&L}

Every day the desk explains its P\&L by \omterm{def:dv:greeks-and-the-hedging-pnl:greeks}{Greeks} and compares the explanation with a full revaluation. With $\delta\sigma_i$ the change of option $i$'s
own mark,
\[
\text{P\&L}\approx\Delta\,\delta S+\tfrac12\Gamma\,\delta S^2+\Theta\,\delta t+\sum_i\bigl(\mathcal V_i\,\delta\sigma_i+\text{vanna}_i\,\delta S\,\delta\sigma_i
+\tfrac12\,\text{volga}_i\,\delta\sigma_i^2\bigr),
\]
and the remainder is the unexplained P\&L. The build computes both, per option and per day. The marking rule enters twice: through the day's
$\delta\sigma_i$, which under \omterm{def:dv:implied-volatility-and-its-surface:sticky}{sticky delta} moves with the spot, and through the marks themselves.

\begin{example}[Three months of a volatility book]\label{ex:dv:trading-volatility:book}
A book long an at-the-money straddle expiring in 0.3 year and short two 90 puts of the same expiry, delta-hedged daily, over 63 days of a
stochastic-volatility path (the share goes from 100 to 98.6 by way of 94.3; the volatility from 20 to 19). Totals, marked \omterm{def:dv:implied-volatility-and-its-surface:sticky}{sticky strike} (\omterm{def:dv:implied-volatility-and-its-surface:sticky}{sticky
delta} in brackets): gamma 2.87 (2.90), \omterm{def:dv:greeks-and-the-hedging-pnl:greeks}{theta} $-1.90$ ($-1.90$), \omterm{def:dv:greeks-and-the-hedging-pnl:greeks}{vega} 1.65 (1.15), \omterm{def:dv:greeks-and-the-hedging-pnl:greeks}{vanna} $-2.24$ ($-1.67$), \omterm{def:dv:greeks-and-the-hedging-pnl:greeks}{volga} $-0.76$ ($-0.62$), unexplained
$-0.16$ ($-0.09$); total $-0.53$ ($-0.22$). The largest daily unexplained is 0.08
(\cref{fig:dv:trading-volatility:explain}).
\end{example}

The two marks agree on gamma and \omterm{def:dv:greeks-and-the-hedging-pnl:greeks}{theta}, which depend on the path, and disagree on \omterm{def:dv:greeks-and-the-hedging-pnl:greeks}{vega} and \omterm{def:dv:greeks-and-the-hedging-pnl:greeks}{vanna}, which depend on how the surface is assumed to
move. They even disagree on the total, because the book is marked differently along the way; the difference disappears only at expiry. A book
whose \omterm{def:dv:greeks-and-the-hedging-pnl:greeks}{vanna} P\&L is as large as its \omterm{def:dv:greeks-and-the-hedging-pnl:greeks}{vega} P\&L is a book whose marking rule is a position.

\begin{omfigure}
\begin{tikzpicture}
\begin{groupplot}[group style={group size=2 by 1,horizontal sep=1.3cm},
  width=6.6cm,height=5.2cm,xlabel={trading day},xmin=0,xmax=63,ymin=-3.5,ymax=3.5,
  tick label style={font=\small},label style={font=\small},grid=major,grid style={gray!25},table/col sep=comma]
\nextgroupplot[ylabel={cumulative P\&L},title={\small sticky strike},
  legend columns=6,legend cell align=left,legend style={font=\footnotesize,at={(1.08,-0.3)},anchor=north,draw=none,fill=none,/tikz/every even column/.append style={column sep=5pt}}]
\addplot[omThm,thick] table[x=day,y=gamma_ss]{figdata/derivatives/25-trading-volatility/explain.csv};
\addplot[omProp,thick] table[x=day,y=theta_ss]{figdata/derivatives/25-trading-volatility/explain.csv};
\addplot[omDef,thick] table[x=day,y=vega_ss]{figdata/derivatives/25-trading-volatility/explain.csv};
\addplot[violet,thick,dashed] table[x=day,y=cross_ss]{figdata/derivatives/25-trading-volatility/explain.csv};
\addplot[gray,thick,dotted] table[x=day,y=unexplained_ss]{figdata/derivatives/25-trading-volatility/explain.csv};
\addplot[black,very thick] table[x=day,y=total_ss]{figdata/derivatives/25-trading-volatility/explain.csv};
\legend{gamma,theta,vega,vanna + volga,unexplained,total}
\nextgroupplot[title={\small sticky delta},yticklabels={}]
\addplot[omThm,thick] table[x=day,y=gamma_sd]{figdata/derivatives/25-trading-volatility/explain.csv};
\addplot[omProp,thick] table[x=day,y=theta_sd]{figdata/derivatives/25-trading-volatility/explain.csv};
\addplot[omDef,thick] table[x=day,y=vega_sd]{figdata/derivatives/25-trading-volatility/explain.csv};
\addplot[violet,thick,dashed] table[x=day,y=cross_sd]{figdata/derivatives/25-trading-volatility/explain.csv};
\addplot[gray,thick,dotted] table[x=day,y=unexplained_sd]{figdata/derivatives/25-trading-volatility/explain.csv};
\addplot[black,very thick] table[x=day,y=total_sd]{figdata/derivatives/25-trading-volatility/explain.csv};
\end{groupplot}
\end{tikzpicture}

\omcaption{Cumulative daily explain of a delta-hedged volatility book (long an at-the-money straddle, short two 90 puts) over 63 days of a
stochastic-volatility path, marked \omterm{def:dv:implied-volatility-and-its-surface:sticky}{sticky strike} (left) and \omterm{def:dv:implied-volatility-and-its-surface:sticky}{sticky delta} (right). Gamma and \omterm{def:dv:greeks-and-the-hedging-pnl:greeks}{theta} agree; \omterm{def:dv:greeks-and-the-hedging-pnl:greeks}{vega} and the cross terms depend on the
marking rule. Data: the tutorial.}\label{fig:dv:trading-volatility:explain}
\end{omfigure}

\section{The variance risk premium in practice}

Chapter 14 defined the \omterm{def:dv:variance-swaps-and-volatility-derivatives:vrp}{variance risk premium} as the gap between implied and expected realised variance. Bakshi and Kapadia measured it through
delta-hedged S\&P~500 options: the hedged positions underperform zero, more so when volatility is high, which is the sign of a negative market
volatility risk premium for the buyer, and of a positive one for the seller. Israelov and Nielsen split a covered call into its equity, short-volatility
and equity-reversal parts. The short-volatility part had a realised Sharpe ratio close to 1.0, though it carried less than a tenth of the strategy's
risk.

\begin{dated}{2026-09}{A published short-volatility benchmark}\label{dat:dv:trading-volatility:put}
Cboe's S\&P~500 PutWrite Index (PUT) tracks a collateralised short-put strategy: at-the-money one-month S\&P~500 puts sold monthly, usually on the
third Friday, over a Treasury-bill account in one- and three-month bills. The strike is the listed strike closest to, but not above, the last
index value reported before 11:00 a.m.\ Eastern time (methodology document, accessed 2026-09).
\end{dated}

The premium is paid for bearing crash risk, and the distribution of a seller's returns shows it: many small gains, a few large losses. On 5 February
2018 the S\&P~500 fell 4\% while the VIX jumped 20 points, the kind of day that takes months of premium.

\begin{example}[Ten years of selling straddles]\label{ex:dv:trading-volatility:vrp}
Every month a desk sells a one-month at-the-money straddle at the volatility its model expects plus two points, an assumed premium, and hedges
daily at that volatility. The share follows a stochastic-volatility path with crashes (on average one a year, $-6\%$). Over 120 months the average
implied volatility sold is 23.5 against 20.6 realised. The P\&L per month averages $+0.44\%$ of the spot, with a standard deviation of 1.33 (a
Sharpe ratio of 1.15 a year), 74\% of months positive, a best month of $+2.89\%$ and a worst of $-5.04\%$; the skewness is $-1.35$
(\cref{fig:dv:trading-volatility:vrp}).
\end{example}

\begin{omfigure}
\begin{tikzpicture}
\begin{axis}[width=10.5cm,height=5cm,ybar,bar width=8pt,
  xlabel={monthly P\&L (\% of spot)},ylabel={months},
  tick label style={font=\small},label style={font=\small},
  xmin=-6,xmax=3.5,ymin=0,ymax=27,grid=major,grid style={gray!25},
  table/col sep=comma]
\addplot[omProp,fill=omProp!40] table[x=centre,y=count]{figdata/derivatives/25-trading-volatility/shortvol.csv};
\end{axis}
\end{tikzpicture}

\omcaption{Monthly P\&L of selling a delta-hedged one-month straddle at two points above the expected volatility, 120 months of a simulated path with
crashes: most months gain, the left tail is long. Data: the tutorial.}\label{fig:dv:trading-volatility:vrp}
\end{omfigure}

\section{Tutorial: scalping and explaining}\label{tut:dv:trading-volatility:tut}

\begin{tutorial}
\textbf{Goal.} Run a delta-hedged straddle along constructed paths and compare the P\&L with the cash-gamma sum; hold a straddle and a calendar on an
unchanged curve; explain a risk reversal's bad day and a volatility book's three months under \omterm{def:dv:implied-volatility-and-its-surface:sticky}{sticky strike} and \omterm{def:dv:implied-volatility-and-its-surface:sticky}{sticky delta}; run a short-volatility
programme for ten years. \textbf{End state:} the four figures, the risk-reversal table and the numbers of the weekend problem.
\begin{tutsteps}
\item \textbf{The explain} of one period, per option at its own mark, against full revaluation:
\omcode{code/firm/volpnl/firm_volpnl.py}{69}{92}{Greek explain against full revaluation.}\label{lst:dv:trading-volatility:explain}
\item \textbf{The scalp}: daily hedges at the implied volatility, and the cash-gamma sum:
\omcode{code/firm/volpnl/firm_volpnl.py}{106}{119}{A hedged straddle along a path.}\label{lst:dv:trading-volatility:scalp}
\item \textbf{Run} \texttt{dv\_volpnl.scalp()}, \texttt{carry()}, \texttt{risk\_reversal()}, \texttt{book\_explain()}, \texttt{short\_vol()} and
\texttt{fig\_volpnl.py}.
\end{tutsteps}
\textbf{What to change next.} Hedge the wrong-days straddle twice a day and at the realised volatility; add a skew move to the book's surface; sell
puts instead of straddles in the programme and compare the tails.
\end{tutorial}

\section{Build: volatility-book P\&L}\label{bld:dv:trading-volatility:volpnl}

\begin{build}
\bfield{Purpose} The miniature firm's volatility-book P\&L: hedged positions revalued in full, and the daily explain by Greek against it, under
either marking rule.
\bfield{Interface} \texttt{Option(strike, expiry, right, qty)}; \texttt{SkewSurface(atm, skew, ref)} with \texttt{vol(strike, tau, spot)}
(\texttt{ref=None}: \omterm{def:dv:implied-volatility-and-its-surface:sticky}{sticky delta}) and \texttt{shifted}; \texttt{value}, \texttt{book\_greeks}; \texttt{explain(book, s0, s1, t0, t1, surf0, surf1,
hedge)}; \texttt{hedged\_pnl(path, book, surfaces)}; \texttt{gamma\_scalp(returns, strike, days, vol)}; \texttt{roll\_down(atm, tau, horizon)}.
\bfield{Rules} Every option is explained at its own mark; the explain is always shown against full revaluation, never instead of it; the marking
rule is an explicit input.
\bfield{Acceptance tests} \texttt{code/firm/volpnl/tests/}: the unexplained part is small for small moves and grows as the cube of the move; the
hedge removes the delta P\&L; sticky-strike marks do not move with the spot and sticky-delta marks do; a straddle hedged along a path realising its
implied volatility makes almost nothing; roll-down on flat and sloping curves.
\bfield{Stretch} \omterm{def:dv:greeks-and-the-hedging-pnl:greeks}{Theta} and roll-down separated in the daily explain; a \omterm{def:dv:greeks-and-the-hedging-pnl:greeks}{vega} explain by expiry bucket (chapter 26); the explain of chapter 18's
\omterm{def:dv:autocallables:autocallable}{autocallables} from their bumped \omterm{def:dv:greeks-and-the-hedging-pnl:greeks}{Greeks}.
\end{build}

\begin{omsources}
\item G.~Bakshi and N.~Kapadia, ``Delta-hedged gains and the negative market volatility risk premium'', \emph{Review of Financial Studies} 16(2)
(2003) 527--566.
\item R.~Israelov and L.~N.~Nielsen, ``Covered calls uncovered'', \emph{Financial Analysts Journal} 71(6) (2015) 44--57.
\item P.~Carr and L.~Wu, ``Variance risk premiums'', \emph{Review of Financial Studies} 22(3) (2009) 1311--1341.
\item Cboe Global Indices, ``Cboe S\&P 500 PutWrite Indices methodology'' (accessed 2026-09).
\item Bank for International Settlements, \emph{Quarterly Review}, March 2018, box on the equity market turbulence of 5 February.
\end{omsources}

\section{Exercises}

\begin{exercise}[$\star$]\label{exo:dv:trading-volatility:1}
A one-month at-the-money straddle on a share at 100 has a \omterm{def:dv:greeks-and-the-hedging-pnl:cashgamma}{cash gamma} of about 770 at inception. What does a day with a 2\% move earn at 18
implied volatility? A day with no move?
\end{exercise}

\begin{exercise}[$\star$]\label{exo:dv:trading-volatility:2}
Why does a long straddle lose money on the wrong-days path although realised volatility beats implied by seven points?
\end{exercise}

\begin{exercise}[$\star$]\label{exo:dv:trading-volatility:3}
Split the three-month straddle's one-month carry into \omterm{def:dv:greeks-and-the-hedging-pnl:greeks}{theta} and roll-down, and say what an inverted curve would change.
\end{exercise}

\begin{exercise}[$\star\star$]\label{exo:dv:trading-volatility:4}
Explain why the calendar earns positive carry and what it pays for it.
\end{exercise}

\begin{exercise}[$\star\star$]\label{exo:dv:trading-volatility:5}
Show that under \omterm{def:dv:implied-volatility-and-its-surface:sticky}{sticky delta}, with the at-the-money level fixed, every strike's volatility changes by $-s\,\ln(S_1/S_0)/\sqrt\tau$ when spot
moves from $S_0$ to $S_1$ (skew $s$), and compute it for the risk reversal's day.
\end{exercise}

\begin{exercise}[$\star\star$]\label{exo:dv:trading-volatility:6}
Why do gamma and \omterm{def:dv:greeks-and-the-hedging-pnl:greeks}{theta} agree under the two marking rules in the book's explain, while \omterm{def:dv:greeks-and-the-hedging-pnl:greeks}{vega} and \omterm{def:dv:greeks-and-the-hedging-pnl:greeks}{vanna} do not?
\end{exercise}

\begin{exercise}[$\star\star\star$]\label{exo:dv:trading-volatility:7}
\emph{Coding.} Rerun the short-volatility programme with no premium. What are the average P\&L and the Sharpe ratio? Is the average
distinguishable from zero, and how can the implied volatility sold exceed the average realised one without a premium?
\end{exercise}

\begin{exercise}[$\star\star\star$]\label{exo:dv:trading-volatility:8}
\emph{Find the flaw.} ``Our short-volatility book has a Sharpe ratio above 1 and three winning months out of four; its risk is small.''
\end{exercise}

\section{Problem: The Wrong Days}

\begin{problem}[{Weekend problem --- when 25 realised loses to 18 implied}]\label{pb:dv:trading-volatility:1}
The trader of the opening bought a 21-day at-the-money straddle on a share at 100 at 18 volatility and hedged it daily. The share realised 25\%.

\medskip\noindent\textbf{Part I --- The trade.}
\begin{enumerate}
\item What did the straddle cost, and what was the trader betting on?
\item Give the daily move that realises 25\% evenly, and the P\&L of the even path.
\item Give the daily \omterm{def:dv:greeks-and-the-hedging-pnl:cashgamma}{cash gamma} at inception, and the breakeven daily move at 18 implied.
\item How much does the even path's P\&L differ from the cash-gamma sum, and why?
\item What would a \omterm{def:dv:variance-swaps-and-volatility-derivatives:vs}{variance swap} on the same notional have paid on any of the three paths?
\end{enumerate}

\medskip\noindent\textbf{Part II --- The wrong days.}
\begin{enumerate}[resume]
\item Describe the wrong-days path and check that it realises 25\%.
\item Split its P\&L between the quiet weeks and the wild week.
\item Where was the spot during the wild week, and what was the straddle's gamma there compared with inception?
\item What does the right-days path make, and why less than the even path?
\item What does the comparison say about the \omterm{def:dv:trading-volatility:spread}{implied--realised spread} as a trading signal?
\end{enumerate}

\medskip\noindent\textbf{Part III --- Other ways to be long volatility.}
\begin{enumerate}[resume]
\item How would re-striking the straddle when the spot drifted have changed the result?
\item What is the carry of holding a three-month straddle for a month on the chapter's curve, and how much of it is roll-down?
\item What does a vega-neutral calendar earn per day, and what risk does it take?
\item On a 5\% fall, what does a short-put, long-call risk reversal make under \omterm{def:dv:implied-volatility-and-its-surface:sticky}{sticky strike} and under \omterm{def:dv:implied-volatility-and-its-surface:sticky}{sticky delta}?
\item Why must a desk revalue such a day in full rather than by \omterm{def:dv:greeks-and-the-hedging-pnl:greeks}{Greeks}?
\end{enumerate}

\medskip\noindent\textbf{Part IV --- Judgement.}
\begin{enumerate}[resume]
\item Was the trader wrong about volatility?
\item What would you report in the P\&L attribution for the month?
\item How does the short-volatility programme's distribution explain why sellers are paid?
\item State the \emph{named result}: the gamma P\&L of a long straddle when the same realised volatility comes in moves concentrated on low-gamma days,
against uniformly spread moves.
\item In one sentence: what does a delta-hedged option pay?
\end{enumerate}
\end{problem}

\section{Interview questions}

\begin{interviewq}[$\star$ \iqroles{trader}]\label{iq:dv:trading-volatility:1}
You are long a delta-hedged straddle. Where does your P\&L come from?
\end{interviewq}

\begin{interviewq}[$\star\star$ \iqroles{trader, researcher}]\label{iq:dv:trading-volatility:2}
Realised volatility came in above implied and your hedged straddle lost money. How?
\end{interviewq}

\begin{interviewq}[$\star\star$ \iqroles{trader}]\label{iq:dv:trading-volatility:3}
What is the carry of a calendar spread on an upward-sloping term structure?
\end{interviewq}

\begin{interviewq}[$\star\star$ \iqroles{risk, developer}]\label{iq:dv:trading-volatility:4}
How would you build a daily P\&L explain for an options book, and what do you do with the unexplained part?
\end{interviewq}

\begin{interviewq}[$\star\star$ \iqroles{trader, risk}]\label{iq:dv:trading-volatility:5}
The same book shows different \omterm{def:dv:greeks-and-the-hedging-pnl:greeks}{vega} P\&L under \omterm{def:dv:implied-volatility-and-its-surface:sticky}{sticky strike} and \omterm{def:dv:implied-volatility-and-its-surface:sticky}{sticky delta}. Which is right?
\end{interviewq}

\begin{interviewq}[$\star\star\star$ \iqroles{researcher}]\label{iq:dv:trading-volatility:6}
Is the \omterm{def:dv:variance-swaps-and-volatility-derivatives:vrp}{variance risk premium} a free lunch? How would you size a short-volatility strategy?
\end{interviewq}
